% \newpage
\section*{Задание 5: Визуализация технической архитектуры}
\subsection{Постановка задачи}
\noindent
Архитектура системы: Процесс разработки (DevOps CI/CD) \\
Ключевые компоненты (узлы): Код (ветка), Сборка (печь), Тест (микроскоп), Релиз (ракета) \\
Визуальная метафора связей: Железнодорожные пути \\
Стиль визуализации: Линейная схема слева направо \\
Цветовое кодирование зон: Градиент от фиолетового (Код) к оранжевому (Продакшн) \\
\newline
Требуется создать изометрическую или плоскую схему, где этапы CI/CD представлены соответствующими метафорическими объектами, соединёнными рельсами, с чётким направлением слева направо и плавным переходом цвета.
\newpage
\subsection{Гипотеза и декомпозиция}
\indent
Для достижения управляемого результата в упрощённом интерфейсе была применена декомпозиция сцены на следующие элементы: \\
Субъекты (узлы): «ветка дерева, символизирующая код», «промышленная печь для сборки», «лабораторный микроскоп для тестирования», «стартующая ракета для релиза» \\
Связи: «железнодорожные пути, соединяющие объекты в линию», «стрелки указывают направление движения» \\
Композиция: «горизонтальная композиция слева направо», «изометрический вид» \\
Цветовая гамма: «градиент от фиолетового (слева) через синий к оранжевому (справа)», «объекты подсвечены соответствующими цветами» \\
Стиль: «чистая векторная иллюстрация», «плоский дизайн с лёгкой 3D-глубиной», «минималистичные детали» \\
Требования к чистоте: «без текста, без подписей, без лишних объектов»
\subsection{Журнал генерации и отладки}
\textbf{Итерация 1} (наивный промпт)\\
Промпт: \textit{<<Схема DevOps CI/CD: слева ветка кода, потом печь сборки, затем микроскоп тестирования, справа ракета релиза. Соедини их железной дорогой. Градиент от фиолетового к оранжевому. Стиль – плоская иллюстрация.»}
\newpage
Результат:
\begin{figure}[H]
	\centering
	\includegraphics[width=16cm]{z5/1a.png}
	\caption{Итерация 1, модель (a)}
\end{figure}
% \newpage
\begin{figure}[H]
	\centering
	\includegraphics[width=16cm]{z5/1b.png}
	\caption{Итерация 1, модель (b)}
\end{figure}
\newpage
\indent \textbf{Анализ ошибок:} \\
Модель восприняла метафоры буквально (ветка дерева вместо абстрактной «ветки кода»), нарушена линейная композиция, объекты имеют несоразмерные масштабы. Отсутствует единый стиль. \\
\newline
\indent \textbf{Корректировка:} \\
Уточнить каждую метафору, добавить слова «абстрактная», «символизирует», явно задать горизонтальную линию, запретить лишние детали. \\
\newline
\indent \textbf{Итерация 2} \\
Промпт: \textit{«Чистая инфографика в плоском стиле. Горизонтальная композиция слева направо. Объекты: слева – абстрактная ветка кода (стилизованная ветка с узлами в виде строк кода), затем – промышленная печь для сборки (цилиндрическая, с паром), затем – лабораторный микроскоп для тестирования, справа – стартующая ракета для релиза. Все объекты соединены прямой железной дорогой (рельсы и шпалы). Цветовой градиент фона от фиолетового (слева) к оранжевому (справа). Без текста, без подписей.»}
\newpage
Результат:
\begin{figure}[H]
	\centering
	\includegraphics[width=16cm]{z5/2a.png}
	\caption{Итерация 2, модель (a)}
\end{figure}
% \newpage
\begin{figure}[H]
	\centering
	\includegraphics[width=16cm]{z5/2b.png}
	\caption{Итерация 2, модель (b)}
\end{figure}
\newpage
\indent \textbf{Анализ ошибок:} \\
Модель исправила композицию, но не смогла удержать единую стилистику «плоская инфографика». Метафоры остались слишком натуралистичными. Градиент интерпретирован как разделение на цветные зоны, а не плавный переход. \\
\newline
\indent \textbf{Корректировка:} \\
Ввести более строгие стилевые ограничения («векторный арт», «силуэты», «минимализм»), явно указать «плавный градиент по всей длине», заменить «ветку кода» на «свёрнутый пергамент с кодом» или «ветвь с шестерёнками». \\
\newline
\indent \textbf{Итерация 3} (финальная) \\
Промпт: \textit{«Изометрическая векторная иллюстрация в стиле чистой инфографики. Горизонтальная схема слева направо. Четыре ключевых объекта, соединённых прямыми железнодорожными путями:
Слева – символ кода: раскрытая книга с текстом программного кода, из которой растёт абстрактная ветвь с маленькими шестерёнками.
Второй объект – промышленная печь для сборки (цилиндрическая камера с элементами индикации тепла).
Третий – лабораторный микроскоп, символизирующий тестирование, с увеличительными линзами.
Справа – стартующая ракета, символизирующая релиз.
Вся сцена в едином плоском стиле, без теней, без реалистичных текстур. Фон – плавный горизонтальный градиент от тёмно-фиолетового (слева) через насыщенный синий к ярко-оранжевому (справа). Цветовая подсветка объектов соответствует зоне градиента. Объекты расположены строго на одной горизонтальной линии, рельсы идут ровно под ними. Без текста, без подписей, без лишних деталей.»}
\newpage
Результат:
\begin{figure}[H]
	\centering
	\includegraphics[width=16cm]{z5/3a.png}
	\caption{Итерация 3, модель (a)}
\end{figure}
% \newpage
\begin{figure}[H]
	\centering
	\includegraphics[width=16cm]{z5/3b.png}
	\caption{Итерация 3, модель (b)}
\end{figure}
\newpage
\indent Получено изображение, полностью удовлетворяющее требованиям: объекты стилизованы в единой векторной манере, имеют чёткие контуры, расположены по горизонтали, соединены прямыми рельсами. Цветовой градиент фона плавный, объекты частично подсвечены цветами градиента. Отсутствуют надписи и посторонние элементы. Метафоры узнаваемы и не противоречат технической тематике.
\subsection{Сравнение моделей}
\begin{table}[h!]
	\centering
	\begin{tabular}{|l|p{5cm}|p{5cm}|}
		\hline
		\textbf{Критерий} & \textbf{модель (a)} & \textbf{модель (b)} \\
		\hline
		Следование промпту & 5 & 5 \\
		\hline
		Качество деталей & 5 & 5 \\
		\hline
		Текст/Схемы & 5 & 5 \\
		\hline
		Цветопередача  & 5 & 5 \\
		\hline
		\textbf{Итог} & 5 & 5\\
		\hline
	\end{tabular}
\end{table}
Обе модели показали хороший результат, также можно отметить что модель (a) пытается упростить, а модель (b) добавляет много мелких деталей. Можно раскрыть потенциал модели (b), написав более точный промпт, чтобы сгенерировать более детальное изображение.
\newpage
\subsection{Финальный промпт}
Промпт: \textit{«Изометрическая векторная иллюстрация в стиле чистой инфографики. Горизонтальная схема слева направо. Четыре ключевых объекта, соединённых прямыми железнодорожными путями:
Слева – символ кода: раскрытая книга с текстом программного кода, из которой растёт абстрактная ветвь с маленькими шестерёнками.
Второй объект – промышленная печь для сборки (цилиндрическая камера с элементами индикации тепла).
Третий – лабораторный микроскоп, символизирующий тестирование, с увеличительными линзами.
Справа – стартующая ракета, символизирующая релиз.
Вся сцена в едином плоском стиле, без теней, без реалистичных текстур. Фон – плавный горизонтальный градиент от тёмно-фиолетового (слева) через насыщенный синий к ярко-оранжевому (справа). Цветовая подсветка объектов соответствует зоне градиента. Объекты расположены строго на одной горизонтальной линии, рельсы идут ровно под ними. Без текста, без подписей, без лишних деталей.»}
\subsection{Вывод}
Ключевые техники, обеспечившие успех:
\begin{itemize}
	\item Декомпозиция метафор: каждому техническому этапу присвоен визуальный образ, избегающий буквального прочтения.
	\item Строгое задание композиции: «горизонтальная схема слева направо», «рельсы ровно под ними» – помогло модели избежать хаотичного размещения.
	\item Управление цветом через контекст: вместо параметра CFG Scale использованы развёрнутые описания градиента и связки цвета объектов с зоной.
	\item Негативные ограничения: «без текста, без подписей, без лишних деталей» – фильтрует артефакты.
\end{itemize}